%---------%---------%---------%---------%---------%---------%---------%
% Upper-case    A B C D E F G H I J K L M N O P Q R S T U V W X Y Z
% Lower-case    a b c d e f g h i j k l m n o p q r s t u v w x y z
% Digits        0 1 2 3 4 5 6 7 8 9
% Exclamation   !           Double quote "          Hash (number) #
% Dollar        $           Percent      %          Ampersand     &
% Acute accent  '           Left paren   (          Right paren   )
% Asterisk      *           Plus         +          Comma         ,
% Minus         -           Point        .          Solidus       /
% Colon         :           Semicolon    ;          Less than     <
% Equals        =           Greater than >          Question mark ?
% At            @           Left bracket [          Backslash     \
% Right bracket ]           Circumflex   ^          Underscore    _
% Grave accent  `           Left brace   {          Vertical bar  |
% Right brace   }           Tilde        ~

\documentclass{entcs}
\usepackage{prentcsmacro}
\usepackage[leqno]{amsmath}
\usepackage{amssymb}
\usepackage{latexsym}
\usepackage{mathrsfs,euscript}


%%%%%%%%%%%%%%%%%%%%%%%%%%%%%%%%%%%%%%%%%%%%%%%%%%%%%%%%%
% Setting the correct page numbers                      %
%%%%%%%%%%%%%%%%%%%%%%%%%%%%%%%%%%%%%%%%%%%%%%%%%%%%%%%%%

  %%%%%%%%%%%%%%%%%%%%%%%%%%%%%%%%%%%%%%%%%%%%%%%%%%%%%
%%%%
%%%%   Calculates and sets the correct page
%%%%   numbers for a particular article.
%%%%
%%%%%%%%%%%%%%%%%%%%%%%%%%%%%%%%%%%%%%%%%%%%%%%%%%%%%
%%%%
%%%%   usage: Insert the following lines before
%%%%          the \begin{document} into the article
%%%%
%%%%            \input procnum
%%%%            \numbering{<database>}{<articleId>}
%%%%
%%%%   Replace <articleId> with the id for the current article
%%%%   and <database> by the name of the textfile containing
%%%%   an entry for each article:
%%%%
%%%%     \article[<optRemark>]{<paperId>}{noPages}{author}{title}
%%%%
%%%%%%%%%%%%%%%%%%%%%%%%%%%%%%%%%%%%%%%%%%%%%%%%%%%%%%%%

\newcounter{startPage}
\setcounter{startPage}{1}

%% Define a command for the entries in the database
%% of articles :
%%     A counter is created from the <articleId>
%%     and initialized to the current startPage
%%     The startPage counter is incremented by the
%%     size of the current article :
%%
\newcommand{\article}[5][]{
  \newcounter{#2}
  \setcounter{#2}{\value{startPage}}
  \addtocounter{startPage}{#3}
}


%% The argument to the following command will be
%% the <articleId>. The corresponding counter's
%% value tells where to start the numbering
%% of this article.

\newcommand{\numbering}[2]{
    \input{#1}
    \setcounter{page}{\value{#2}}
}

%%%%%%%%%%%%%%%%%%%%%%%%%%%%%%%%%%%%%%%%%%%%%%%%%%%%%

  \numbering{cmcs03db}{adamek}

%%%%%%%%%%%%%%%%%%%%%%%%%%%%%%%%%%%%%%%%%%%%%%%%%%%%%%%%%



%%%%%%%%%%%%%%%%%%%%%%%%%%%%%%%%%%%%%%%%%%%%%%%%%%%%%%%%%
% CALLING XY-Pic version 3.5                            %
%%%%%%%%%%%%%%%%%%%%%%%%%%%%%%%%%%%%%%%%%%%%%%%%%%%%%%%%%
\usepackage[matrix,arrow]{xy}


\newcommand{\Set}{\ensuremath{\St}}
\newcommand{\aps}{\ensuremath{\mathrel{\approx^*}}}
\newcommand{\sims}{\ensuremath{\mathrel{\sim^*}}}
\newcommand{\simx}{\ensuremath{\mathrel{\sim_X}}}
\newcommand{\simxs}{\ensuremath{\mathrel{\sim_X^*}}}
\newcommand{\vxs}{\ensuremath{\mathrel{\approx_X^*}}}
\newcommand{\Pf}{\ensuremath{\mathscr{P}_{\!\!f}}}
\newcommand{\AAA}{\ensuremath{\mathscr{A}}}
\newcommand{\PP}{\ensuremath{\mathscr{P}}}
\newcommand{\TT}{\ensuremath{\mathscr{T}}}
\newcommand{\RR}{\ensuremath{\mathscr{R}}}
\newcommand{\SSS}{\ensuremath{\mathscr{S}}}
\newcommand{\llll}{l}
\newcommand{\dia}{\mathbin{\Diamond}}
\newcommand{\ctv}{\mathbin{\Box}}
\newcommand{\ac}{\text{\textnormal{ac}}}

\DeclareMathOperator{\St}{\textnormal{\textbf{Set}}}
\DeclareMathOperator{\id}{id}
\DeclareMathOperator{\Id}{Id}
\DeclareMathOperator{\Alg}{\textnormal{\textbf{Alg}}}
\DeclareMathOperator{\Coalg}{\textnormal{\textbf{Coalg}}}
\DeclareMathOperator*{\colim}{colim}
\DeclareMathOperator{\inl}{inl}
\DeclareMathOperator{\inr}{inr}
\DeclareMathOperator{\CIM}{CIM}
\DeclareMathOperator{\IM}{IM}
\DeclareMathOperator{\card}{card}

\newenvironment{pv}[1]{\begin{trivlist}
\item[\hspace{\labelsep}\normalfont\bfseries#1]}%
{\end{trivlist}}

\def\lastname{Ad\'{a}mek}

\begin{document}

\begin{frontmatter}
\title{On a Description of Terminal Coalgebras and Iterative
Theories}
\author{Ji\v r\'\i\ Ad\'amek\thanksref{mail}\thanksref{CZ}}
\address{Technical University of Braunschweig}
\thanks[mail]{Email:\href{mailto:adamek@iti.cs.tu-bs.de}
        {\texttt{\normalshape adamek@iti.cs.tu-bs.de}}}
\thanks[CZ]{The support of the Grant Agency of the Czech
            Republic under the Grant No.~201/02/0148
            is gratefully acknowledged.}
\begin{keyword}
terminal coalgebra, iterative theory, finitary functor, basic equation
\end{keyword}
\begin{abstract}
A concrete description of terminal coalgebras, $T$, of all finitary
endofunctors of \Set\ is presented: each such a functor is a quotient of
the polynomial endofunctor of some finitary signature~$\Sigma$
modulo some ``basic'' equations. Then $T$~can be described as the
algebra of all infinite $\Sigma$-labeled trees modulo the congruence
obtained by applying the basic equations finitely or infinitely
many times (a concept defined below). As a consequence, free
iterative theories in the sense of Calvin Elgot are described over
all finitary endofunctors of \Set: they are the theories of all
rational $\Sigma$-labeled trees (i.e., trees having only finitely
many subtrees) modulo the above congruence.
\end{abstract}
\end{frontmatter}


\section{Introduction}
Iterative algebraic theories have been introduced by Calvin Elgot
(see~\cite{e}) in his attempt to model potentially infinite
computations algebraically. For example, a computation using binary
operations $\alpha$ and~$\beta$ and having the following form
\begin{equation}
\alpha(0, \beta(1,\alpha(0, \beta(1,\dots
\end{equation}
can be described by iterative equations as follows:
\begin{equation}
x_1=
\vcenter{
\xy
\POS (05,10)*+{\alpha} = "t"
   , (00,00)*+{0}      = "l"
   , (10,00)*+{x_2}    = "r"
%%%
\POS "l" \ar@{-} "t"
\POS "r" \ar@{-} "t"
\endxy
}
\quad
\quad
x_2=
\vcenter{
\xy
\POS (05,10)*+{\beta}  = "t"
   , (00,00)*+{1}      = "l"
   , (10,00)*+{x_1}    = "r"
%%%
\POS "l" \ar@{-} "t"
\POS "r" \ar@{-} "t"
\endxy
}
\end{equation}
One of the main result of Elgot's group is a description of free
iterative theories generated by an arbitrary signature~$\Sigma$ (see
~\cite{ebt}): it is the theory of all rational
$\Sigma$-trees\footnote{Trees are always considered up-to
isomorphism (of labeled trees) throughout the paper.}, i.e.,
$\Sigma$-trees which have only finitely many subtrees. Example: the
tree~(1) above is rational, it has 4~subtrees.

One fundamental lack of Elgot's theory is that almost no other
concrete examples have been presented. In the
present paper we take the first steps towards new examples: we
describe all iterative theories freely generated by finitary
endofunctors of \Set. Or, equivalently, freely generated by
equational presentations using \textit{basic equations} only, where
a basic equation is an equation between two \textit{flat terms}
which are terms $\sigma(x_1,\dots,x_n)$ where $\sigma\in\Sigma$ is
an $n$-ary operation symbol and $x_1,\dots,x_n$ are (not necessarily
distinct) variables.

\begin{pv}{1.1 Example}
Let $H\colon\Set\to\Set$
be the functor assigning
to every set~$X$ the set~$HX$ of all nonordered
pairs~$\{x,y\}$ in~$X$. This is a quotient of the functor
$H_\Sigma\colon X\mapsto X\times X$ via the epitransformation
$\varepsilon\colon H_\Sigma\to H$ defined
by $\varepsilon(x,y)=\{x,y\}$. Or,
equivalently, we consider the signature~$\Sigma$ of one binary
operation (say,~$\diamond$) satisfying the commutativity law.
Observe that the equation
\[x\diamond y=y\diamond x\]
is indeed basic. The free iterative theory is formed by all
rational binary trees modulo the congruence \aps\ merging two
rational trees $s$ and~$r$ iff we can obtain~$r$ from~$s$ by
applying the commutativity law (i.e., by swapping the two children
of some parent) finitely or infinitely many times.

This ``infinite application'' of equations can be described simply
and precisely by introducing the operator
\[\partial_k\]
which cuts every tree at depth~$k$ and labels all leaves of
depth~$k$ by an auxiliary symbol~$\perp$. Let us denote by~$\sim$
the congruence defining a free algebra of our variety. That is, if
we think of terms as finite binary trees, then $s\sim r$ means that
we can obtain~$r$ from~$s$ by applying the commutativity law
finitely many times. Then for all rational trees $s$ and~$r$ we can
define
\[s\aps r\qquad\text{iff}\qquad \partial_ks\sim\partial_kr\quad
\text{for all $k\in\omega$.}\]
\end{pv}

\begin{pv}{Remark}
The above definition of~\aps\ is meaningful for all basic equational
presentations, since the congruence~$\sim$ defining the free algebra
of the given variety satisfies
\[s\sim r\qquad\text{implies}\qquad \partial_ks\sim\partial_kr\quad
\text{for all $k\in\omega$}\]
(for all finite trees $s$ and~$r$). In contrast, this is not
meaningful e.g. for the associativity law.
\end{pv}

\begin{pv}{1.2 Example}
The iterative theory generated by the finite-power-set functor
\[\Pf X=\{M\subseteq X; M\text{ finite}\}.\]
On morphisms $h\colon X\to Y$, we have
\[\Pf h\colon M\mapsto f[M]\qquad\text{for all $M\in\Pf X$.}\]
There is a natural presentation of~\Pf\ as the quotient functor of
the polynomial functor
\[H_\Sigma X=1+X+(X\times X)+(X\times X\times X)+\dotsb\]
corresponding to the signature~$\Sigma$ which has for every
$n\in\omega$ precisely one $n$-ary operation~$\sigma_n$: define a
natural epitransformation
\[\varepsilon\colon H_\Sigma\to\Pf\]
by
\[(x_1,\dots,x_n)\mapsto\{x_1,\dots,x_n\}.\]
Or, equivalently, consider the following basic equations
\[\sigma_n(y_1,\dots,y_n)=\sigma_m(z_1,\dots,z_m)\]
where $n,m\in\omega$ are arbitrary and $y_i$ and~$z_j$ are arbitrary
variables such that the sets $\{y_1,\dots,y_n\}$
and~$\{z_1,\dots,z_m\}$ are equal.

Here, again, the corresponding iterative theory is formed by all
rational $\Sigma$-trees modulo the congruence~\aps\ obtained by
applying the given equations finitely or infinitely many times.

Let us recall that an initial \Pf-algebra $I$ can be described as
the algebra of all sets of
the class $V_\omega=\bigcup_{n\in\omega}V_n$ of
the cumulative hierarchy (with $V_0=\emptyset$ and $V_{n+1}=\Pf
V_n$). Or directly: $I$~is the algebra of all finite, non-ordered
\textit{extensional trees}, i.e., trees such that any two
distinct siblings
define nonisomorphic subtrees.
%In fact, there is a canonical bijection
%between $V_\omega$ and such trees, as the following examples
%demonstrate:
%\[\psset{radius=2pt,levelsep=5mm,linewidth=.6pt}
%\begin{array}{lcc}
%\text{Set}&\rule{3em}{0em}&\text{tree}\\
%\hline
%\phantom{\{}\emptyset&\rule{0pt}{13pt}&
%\pstree{\TC*}{}\\
%\{\emptyset\}&&
%\raisebox{3mm}{\pstree{\TC*}{\TC*}}\\[13pt]
%\{\emptyset,\{\emptyset\}\}&&
%\raisebox{6mm}{\pstree{\TC*}{\TC* \pstree{\TC*}{\TC*}}}\\[13pt]
%\{\{\emptyset\}\}&&
%\raisebox{6mm}{\pstree{\TC*}{\pstree{\TC*}{\TC*}}}\\
%\quad\vdots&&\vdots
%\end{array}\]
In the latter description we obtain~$I$ as a quotient of the initial
$\Sigma$-algebra~$I_\Sigma$ (of all finite trees---here we can drop
the labeling since $\Sigma$~has a unique operation for any arity)
modulo the following congruence~$\sim$:
\[s\sim r\quad\text{iff $s$ and $r$ have the same extensional
quotient.}\]
(The \textit{extensional quotient} of a finite tree is the extensional tree
obtained by repeated identification of siblings defining isomorphic
subtrees.)
\end{pv}

The above two examples are typical: we will show that every finitary
endofunctor~$H$ of~\Set\ can be presented, for some signature
$\Sigma=\bigcup_{n\in\omega}\Sigma_n$, as a quotient of the
\textit{polynomial functor}~$H_\Sigma$ defined on objects by
\[H_\Sigma X=\coprod_{n\in\omega}\Sigma_n\times X^n\]
and analogously on morphisms, $H_\Sigma f=\coprod_{n\in\omega}
(\id_{\Sigma_n}\times f^n)$. That is, there is an epitransformation
(a natural transformation whose components are epimorphisms)
\[\varepsilon\colon H_\Sigma\to H.\]
Or, equivalently, $H$~represents the equational presentation
by all the basic
equations $\sigma(x_1,\dots,x_n)=\varrho(y_1,\dots,y_m)$ for
all $\sigma\in\Sigma_n$, $\varrho\in\Sigma_m$ and all
variables~$x_i,y_j$ in a set~$X$ such that $\varepsilon_X$~merges
the two elements $\sigma(x_i)$ and~$\varrho(y_j)$ of~$H_\Sigma X$.
Now let $\sim$~be the congruence on the initial $\Sigma$-algebra (of
all finite $\Sigma$-trees) defining a free $H$-algebra, i.e., the
congruence of applications of the given basic equations finitely
many times. Then for the congruence~\sims\ on the algebra of all
finite and infinite trees given by possibly infinite applications of
the given equations we prove that
\begin{enumerate}
\item[(a)]
a terminal $H$-coalgebra is the quotient of the terminal
$\Sigma$-algebra (of all infinite $\Sigma$-trees) modulo~\sims
\end{enumerate}
and
\begin{enumerate}
\item[(b)]
a free iterative theory on~$H$ is the theory of all rational
$\Sigma$-trees modulo the congruence~\aps\ defined analogously
to~\sims\ but for trees with variables.
\end{enumerate}

The method of proving this result is coalgebraic. We first describe
a terminal $H$-coalgebra, for every finitary endofunctor $H$,
as~$T_\Sigma/_{\sims}$
where $T_\Sigma$ is the coalgebra of all (infinite)
$\Sigma$-trees and \sims~is the above congruence. We apply this,
then, to the finitary endofunctor $H({-})+X$. Here
we use the result of~\cite{aamv} that
terminal coalgebras, $TX$, of the functors
$H({-})+X$ form (the object part of) a monad~\TT\
which is completely iterative. In fact, \TT~is a free completely
iterative monad on~$H$. And as proved in~\cite{amv},
a free iterative monad on~$H$ is a submonad of the
monad~\TT\ formed by solutions of all finite systems of flat
equations.

\begin{pv}{Related Results}
I have been much inspired by the description of a terminal
coalgebra of~\Pf\ by M.~Barr~\cite{b}, see Example~1.2 above.
Results of~\cite{aamv} mentioned above
have been independently discovered by L.~Moss~\cite{m}.
\end{pv}

\section{A Description of Terminal Coalgebras}
\begin{pv}{2.1 Assumption}
We assume throughout the present section that a \textit{finitary}
endofunctor~$H$ of~\Set\
is given. Recall that this means one of the following equivalent
properties (see 4.3 in~\cite{at}):
\begin{enumerate}
\item[(i)]
$H$~preserves filtered colimits,
\item[(ii)]
every element of~$HX$, where $X$~is an arbitrary set, lies in the
image of~$Hm$
for some finite subset $m\colon M\hookrightarrow X$
\end{enumerate}
and
\begin{enumerate}
\item[(iii)]
$H$~is a quotient of some polynomial functor~$H_\Sigma$.
\end{enumerate}
The last condition presents~$H$ via a finitary signature~$\Sigma$
and a natural transformation
\[\varepsilon\colon H_\Sigma\to H\]
having epimorphic components. Therefore each component
\[\varepsilon_X\colon \coprod_{n\in\omega}
\Sigma_n\times X^n\to HX\]
is fully described by its kernel equivalence, which we can present
in the form of equations
\[\sigma(x_i)=\varrho(y_j)\]
for $\sigma,\varrho\in\Sigma$ and for tuples
$\sigma(x_i),\varrho(y_j)$ in~$H_\Sigma X$ (where $X$~is a set of
variables including all $x_i$ and $y_j$ and $\sigma(x_i)$ denotes the
$n$-tuple $(x_i)$ in the summand $\{\sigma\}\times X^n$) satisfying
$\varepsilon_X\bigl(\sigma(x_i)\bigr)=
\varepsilon_X\bigl(\varrho(y_j)\bigr)$.
We call the above equations \textit{$\varepsilon$-equations} for short.
Observe that every $\varepsilon$-equation is \textit{basic}, i.e.,
it
equates two \textit{flat}
trees (which means trees of depth one with variables on all leaves).
Conversely,
every variety of $\Sigma$-algebras presented by basic equations is
the category of $H$-algebras for some finitary endofunctor~$H$
of~\Set, see~3.3 in~\cite{at}.
\end{pv}

\begin{pv}{2.2 Examples}
(i) The functor $\PP_{\!2}$ assigning to a set~$A$ the
set~$\PP_{\!2}A$ of all subsets of power at most~2 is a quotient
of~$H_\Sigma$ where
\[\Sigma_0=\{b\}\quad\text{and}\quad \Sigma_2=\{\beta\}.\]
Here
\[\varepsilon_X\colon X\times X+1\to\PP_{\!2}X\]
sends a pair~$(x,y)$ to~$\{x,y\}$ and the unique element of~1
to~$\emptyset$. The $\varepsilon$-equations are all generated by the
commutativity of~$\beta$:
\[\beta(x,y)=\beta(y,x).\]

(ii) Consider the finite-power-set functor~\Pf\ of~1.2. Here we can
use the signature~$\Sigma$ where $\Sigma_n$~contains a unique
$n$-ary operation for any $n=0,1,2,\dots$ and obtain a natural
epitransformation
\[\varepsilon_X\colon1+X+X^2+X^3+\dots\to\Pf X\]
sending an $n$-tuple to the set of its members.
\end{pv}

\begin{pv}{2.3 Notation}
$\Alg H$ denotes the category of \textit{$H$-algebras}, i.e.,
sets~$A$ equipped with a structure morphism $\alpha\colon HA\to A$,
and \textit{homomorphisms}~$f$ between $H$-al\-ge\-bras defined by the
commutativity of the following square\vspace{1mm}
\[                        %% o5
\xy
\xymatrix{
HA
\ar[0,1]^-{\alpha}
\ar[1,0]_{Hf}
&
A
\ar[1,0]^{f}
\\
HA'
\ar[0,1]_-{\alpha'}
&
A'
}
\endxy
\]
\vspace{1mm}

\noindent
An \textit{initial algebra}, i.e., an initial object of~$\Alg H$
(which exists, in fact, is given by~$\colim_{n\in\omega} H^n\emptyset$,
since
$H$~is finitary (see [ A]) is denoted by~$I$ and its structure
morphism by
\[
\xy
\xymatrix@1{
HI
\ar[0,1]^-{\varphi}
&
I.
}
\endxy
\]
\end{pv}

\begin{pv}{2.4 Examples}
(i) If $H=H_\Sigma$, then $\Alg H_\Sigma$~is the usual
category of $\Sigma$-al\-ge\-bras and homomorphisms.
There are several
useful descriptions of~$I_\Sigma$, an initial $\Sigma$-al\-ge\-bra;
here we consider the following:
\[I_\Sigma=\text{the algebra of all finite $\Sigma$-trees.}\]
Recall that a \textit{$\Sigma$-tree} is a (rooted, ordered) labeled
tree such that every node of $n$~children is labeled by an $n$-ary
symbol in~$\Sigma$. The structure morphism
\[\varphi_\Sigma\colon H_\Sigma I_\Sigma\to I_\Sigma\]
is given by tree-tupling.

(ii) An initial $\PP_{\!2}$-algebra~$I$ is just an initial algebra
of the variety of $\Sigma$-al\-ge\-bras [see 2.2~(i)], given by
commutativity of~$\beta$. Thus
\[I=I_{\Sigma}/_\sim\]
where two finite binary trees $s,r\in I_{\Sigma}$ are congruent
iff we can get~$s$ from~$r$ by swapping two siblings
finitely many times.
\end{pv}

\begin{pv}{2.5 Remark}
We have an embedding
\[\Alg H\hookrightarrow\Alg H_\Sigma\]
assigning to every $H$-algebra,
$HA\xrightarrow{\rule{1mm}{0mm}\alpha\rule{1mm}{0mm}}A$, the
$\Sigma$-algebra
\[
\xy
\xymatrix@1{
H_\Sigma A
\ar[0,1]^-{\varepsilon_A}
&
HA
\ar[0,1]^-{\alpha}
&
A.
}
\endxy
\]
In fact, $\Alg H$~is a variety of $\Sigma$-algebras (presented by
the given basic equations). Conversely, every $\Sigma$-algebra~$B$
has a \textit{reflection} in~$\Alg H$: it is the smallest quotient
$j\colon B\to B/_\sim$ such that $B/_\sim$~is an $H$-algebra (more
precisely, a $\Sigma$-algebra whose structure morphism factorizes
trough~$\varepsilon_{B/_\sim}$):\vspace{3mm}
\[
\xy
\xymatrix{
H_\Sigma B
\ar[0,2]^-{\beta}
\ar[1,0]_{H_\Sigma j}
&
&
B
\ar[1,0]^{j}
\\
H_\Sigma(B/_\sim)
\ar[0,1]_-{\varepsilon_{B/_\sim}}
&
H(B/_\sim)
\ar[0,1]_-{\overline{\beta}}
&
B/_\sim
}
\endxy
\]\vspace{1mm}
\end{pv}

\begin{pv}{2.6 Notation}
Given $\varepsilon\colon H_\Sigma\to H$, denote by
\[\Sigma^\perp\]
the signature obtained from~$\Sigma$ by adding a new nullary
symbol~$\perp$. Then $H_\Sigma$~is a subfunctor of~$H_{\Sigma^\perp}$
and we denote by
\[u\colon H_\Sigma\to H_{\Sigma^\perp}\]
the natural transformation formed by the inclusion maps. This allows
us to view every $\Sigma^\perp$-algebra,
$H_{\Sigma^\perp}A\xrightarrow{\rule{1mm}{0mm}\alpha\rule{1mm}{0mm}}A$,
as a $\Sigma$-algebra via
\[
\xy
\xymatrix@1{
H_\Sigma A
\ar[0,1]^-{u_A}
&
H_{\Sigma^\perp}A
\ar[0,1]^-{\alpha}
&
A.
}
\endxy
\]
In particular, the algebra~$I_{\Sigma^\perp}$ of finite
$\Sigma^\perp$-trees is a $\Sigma$-algebra.
\end{pv}

\begin{pv}{Definition}
\itshape
Given a quotient $\varepsilon\colon H_\Sigma\to H$ we denote by
\[J=I_{\Sigma^\perp/_\sim}\]
the reflection of the $\Sigma$-algebra~$I_{\Sigma^\perp}$
in~$\Alg H$\ with the structure morphism~$\varphi$:\vspace{3.5mm}
\[
\xy
\xymatrix{
H_\Sigma I_{\Sigma^\perp}
\ar[0,1]^-{u_{I_{\Sigma^\perp}}}
\ar[1,0]_{H_\Sigma j}
&
H_{\Sigma^\perp}I_{\Sigma^\perp}
\ar[0,1]^-{\varphi_{\Sigma^\perp}}
&
I_{\Sigma^\perp}
\ar[1,0]^{j}
\\
H_\Sigma J
\ar[0,1]_-{\varepsilon_J}
&
HJ
\ar[0,1]_-{\varphi}
&
J
}
\endxy
\]\vspace{.5mm}

\noindent
In other words, for finite $\Sigma^\perp$-trees $s$ and~$r$ we have
\[s\sim r\quad \text{iff $s$ can be obtained from $r$ by using
$\varepsilon$-equations (finitely many times).}\]
\end{pv}

\begin{pv}{2.7 Example}
For $\varepsilon\colon H_\Sigma\to\PP_{\!2}$ as in 2.2 (i) the
algebra~$I_{\Sigma^\perp}$ consists of all finite binary trees with
leaves labeled by $b$ or~$\perp$. And $\sim$~is the congruence of
commutativity of the binary operation, i.e., $s\sim r$ if $s$~can be
obtained from~$r$ by swapping the children of one father finitely
many times.
\end{pv}

\begin{pv}{2.8 Notation}
We will denote by
$\Coalg H$ the category of \textit{$H$-coalgebras}, i.e.,
sets~$A$
equipped with a structure morphism $\alpha\colon A\to HA$,
and \textit{homomorphisms}~$f$ between $H$-coalgebras defined by the
commutativity by the following square\vspace{1mm}
\[
\xy
\xymatrix{
A
\ar[0,1]^-{\alpha}
\ar[1,0]_{f}
&
HA
\ar[1,0]^{Hf}
\\
A'
\ar[0,1]_-{\alpha'}
&
HA'
}
\endxy
\]
\vspace{.1mm}

\noindent
Again, a \textit{terminal coalgebra}, i.e., a terminal object
of~$\Coalg H$, exists due to the finitarity of~$H$ (see~\cite{b}), and
we denote it by~$T$ with the structure morphism
\vspace{1mm}
\[
\xy
\xymatrix@1{
T
\ar[0,1]^-{\psi}
&
HT.
}
\endxy
\]
Recall that by Lambek's Lemma~\cite{l}, $\psi$~is an isomorphism;
thus $T$~can also be viewed as an $H$-algebra.
\end{pv}

\begin{pv}{Example}
A terminal $H_\Sigma$-coalgebra can be described as the coalgebra
\[T_\Sigma\]
of all (finite and infinite) $\Sigma$-trees. Here
\[\psi_\Sigma\colon T_\Sigma\to H_\Sigma T_\Sigma\]
is given by parsing, i.e., the inverse of the tree-tupling.
\end{pv}

\begin{pv}{2.9 Definition}
\itshape
For every $k=0,1,2,\dots$ denote by
\[\partial_k\colon T_\Sigma\to T_{\Sigma^\perp}\]
the function cutting a given tree and level~$k$ at labeling all new
leaves by~$\perp$. Define a congruence~\sims\ on the
$\Sigma$-algebra $H_\Sigma T_\Sigma
\xrightarrow{\rule{1mm}{0mm}\psi_\Sigma^{-1}\rule{1mm}{0mm}}
T_\Sigma$ as follows:
\[s\sims r\quad\text{iff $\partial_ks\sim\partial_kr$ for all
$k=0,1,2,\dotsc$,}\]
where $\sim$~is the congruence of Definition~2.6.
\end{pv}

\begin{pv}{Remark}
We prove in Theorem 2.11 below that \sims~is a congruence for~$H$,
and that the quotient coalgebra
\[J^*=T_\Sigma/_{\sims}\]
of the  $H$-algebra\vspace{1mm}
\[
\xy
\xymatrix{
T_\Sigma
\ar[0,1]^-{\psi_\Sigma}
&
H_\Sigma T_\Sigma
\ar[0,1]^-{\varepsilon_{T_\Sigma}}
&
HT_\Sigma
}
\endxy
\]
is a terminal. We denote by $j^*\colon T_\Sigma\to J^*$ the quotient
map.
\end{pv}

\begin{pv}{2.10 Examples}
(i) For $\varepsilon\colon H_\Sigma\to\PP_{\!2}$ of Example
2.2~(i) we have
\[                                   %% o6
\vcenter{
\xy
\POS (00, 00) *+{b} = "b1"
   , (05, 10)       = "t1"
   , (10, 00)       = "t2"
   , (05,-10) *+{b} = "b2"
   , (15,-10)       = "t3"
   , (10,-20) *+{b} = "b3"
   , (20,-20)       = "t4"
   , (25,-30)        = "x"
%%%
\POS "b1" \ar@{-} "t1"
\POS "t2" \ar@{-} "t1"
\POS "b2" \ar@{-} "t2"
\POS "t3" \ar@{-} "t2"
\POS "b3" \ar@{-} "t3"
\POS "t4" \ar@{-} "t3"
\POS "x" \ar@{.} "t4"
\endxy
}
\quad\sims\quad
\vcenter{
\xy
\POS ( 00, 00) *+{b} = "b1"
   , (-05, 10)       = "t1"
   , (-10, 00)       = "t2"
   , (-05,-10) *+{b} = "b2"
   , (-15,-10)       = "t3"
   , (-10,-20) *+{b} = "b3"
   , (-20,-20)       = "t4"
   , (-25,-30)        = "x"
%%%
\POS "b1" \ar@{-} "t1"
\POS "t2" \ar@{-} "t1"
\POS "b2" \ar@{-} "t2"
\POS "t3" \ar@{-} "t2"
\POS "b3" \ar@{-} "t3"
\POS "t4" \ar@{-} "t3"
\POS "x" \ar@{.} "t4"
\endxy
}
\]
because in~$I_{\Sigma^\perp}$ we clearly have
\[                                   %% o7
\begin{array}{rcrcl}
(k=0)
&\qquad&
\perp
&\sim&
\perp
\\[2mm]
(k=1)
&&
\xy
\POS (00,00) *+{b}     = "b"
   , (10,00) *+{\perp} = "p"
   , (05,10)           = "t"
%%%
\POS "b" \ar@{-} "t"
\POS "p" \ar@{-} "t"
\endxy
&\sim&
\xy
\POS (00,00) *+{\perp} = "b"
   , (10,00) *+{b}     = "p"
   , (05,10)           = "t"
%%%
\POS "b" \ar@{-} "t"
\POS "p" \ar@{-} "t"
\endxy
\\[8mm]
(k=2)
&&
\xy
\POS (00, 00) *+{b} = "b1"
   , (05, 10)       = "t1"
   , (10, 00)       = "t2"
   , (05,-10) *+{b} = "b2"
   , (15,-10) *+{\perp} = "t3"
%%%
\POS "b1" \ar@{-} "t1"
\POS "t2" \ar@{-} "t1"
\POS "b2" \ar@{-} "t2"
\POS "t3" \ar@{-} "t2"
\endxy
&\sim&
\xy
\POS ( 00, 00) *+{b} = "b1"
   , (-05, 10)       = "t1"
   , (-10, 00)       = "t2"
   , (-05,-10) *+{b} = "b2"
   , (-15,-10) *+{\perp} = "t3"
%%%
\POS "b1" \ar@{-} "t1"
\POS "t2" \ar@{-} "t1"
\POS "b2" \ar@{-} "t2"
\POS "t3" \ar@{-} "t2"
\endxy
\end{array}
\]
etc.

(ii) For $\varepsilon\colon H_\Sigma\to\Pf$ of Example 2.2~(ii)
we can describe
\[\parbox{91mm}{$I_{\Sigma^{\phantom{\perp}}}$
as the algebra of all finite
(nonlabeled) trees,\\
$I_{\Sigma^\perp}$ as the algebra of all finite trees with leaves
partially $\phantom{I_{\Sigma^\perp}}$~labeled by~$\perp$,}\]
and
\[\parbox{91mm}{$T_\Sigma$ as the algebra of all finite-branching
trees.}\]
(The operation labels can be left out here since $n$-ary operations
are unique for all~$n$.) Recall the concept of extensional quotient
of Example~1.2. The congruence~$\sim$ on~$I_{\Sigma^\perp}$ is
easily seen to be the following:
\[\text{$r\sim s$ iff the trees $r,s$ have the same extensional
quotient.}\]
It was M.~Barr who used in~\cite{b} the operation~$\partial_k$ to
describe a terminal coalgebra of~\Pf\ as~$T_\Sigma/{\sims}$ for the
above congruence~\sims.
%congruent trees:
%\[                                   %% o8
%\psset{levelsep=7mm,treesep=6mm,radius=2pt}
%\raisebox{11mm}{%
%\pstree{\TC*}{\TC*
%\pstree{\TC*}{\TC*}
%\pstree{\TC*}{\pstree{\TC*}{\TC*}}
%\pstree{\TC*~[tnpos=r,tnsep=6pt]{\raisebox{3pt}{$\dots$}}}{%
%\pstree{\TC*}{\pstree{\TC*}{\TC*}}}}}
%\quad\sim^*\quad
%\raisebox{11mm}{%
%\pstree{\TC*}{%
%\pstree{\TC*}{
%\pstree{\TC*}{
%\pstree{\TC*}{
%\pstree{\TC*}{\TC*~[tnpos=b,tnsep=0pt]{\vdots}}}}}
%\TC*
%\pstree{\TC*}{\TC*}
%\pstree{\TC*}{\pstree{\TC*}{\TC*}}
%\pstree{\TC*~[tnpos=r,tnsep=6pt]{\raisebox{3pt}{$\dots$}}}{%
%\pstree{\TC*}{\pstree{\TC*}{\TC*}}}}}
%\]
\end{pv}

\begin{pv}{2.11 Theorem}
\itshape
For every finitary endofunctor~$H$ a terminal coalgebra can be
described as the quotient
\[T_\Sigma/_{\sims}\]
of a terminal $H_\Sigma$-coalgebra~$T_\Sigma$ (of all
$\Sigma$-trees) modulo the congruence~\sims of Definition 2.9.
\end{pv}

\begin{pv}{Remark}
More detailed: given a presentation of~$H$ as a quotient
$\varepsilon\colon H_\Sigma\to H$
and defining the quotient map $j^*\colon
T_\Sigma\to T_\Sigma/_{\sims}$ as in 2.9, there exists a
unique structure of a coalgebra, $\psi^*$, on~$T_\Sigma/_{\sims}$
making~$j^*$ a homomorphism:\vspace{3.5mm}
\[
\xy
\xymatrix{
T_\Sigma
\ar[0,1]^-{\psi_\Sigma}
\ar[1,0]_{j^*}
&
H_\Sigma T_\Sigma
\ar[0,1]^-{\varepsilon_{T_\Sigma}}
&
HT_\Sigma
\ar[1,0]^{Hj^*}
\\
T_\Sigma/_{\sims}
\ar[0,2]_-{\psi^*}
&
&
H(T_\Sigma/_{\sims})
}
\endxy
\]\vspace{1mm}

\noindent
And the coalgebra~$T_\Sigma/_{\sims}$ is terminal for~$H$.
\end{pv}

We denote $J^*=T_\Sigma/_{\sims}$.

\begin{pf}
We start with some preliminary facts.

(1) For every natural number~$k$ there is unique
morphism~$\overline{\partial}_k$ such that the following square
\vspace{1mm}
\[
\xy
\xymatrix{
T_\Sigma
\ar[0,1]^-{\partial_k}
\ar[1,0]_{j^*}
&
I_{\Sigma^\perp}
\ar[1,0]^{j}
\\
J^*
\ar[0,1]_-{\overline{\partial}_k}
&
J
}
\endxy
\]
\vspace{1mm}

\noindent
commutes. Moreover, the cone $(\overline{\partial})_{k\in\omega}$ has
the property that for every finite set $m\colon M\hookrightarrow
J^*$ there is a~$k$ such that $\overline{\partial}_k{\cdot}m$~is a
monomorphism.

In fact, the first statement follows from the observation that given
$s,r\in T_\Sigma$ with $j^*(s)=j^*(r)$, then
$j(\partial_ks)=j(\partial_kr)$, due to the definition of~\sims.

For the latter statement, find a finite set $m_0\colon
M_0\hookrightarrow T_\Sigma$ such that $M$~is the image of~$M_0$
under~$j^*$ and $j^*m_0$~is a monomorphism. For this finite
set~$M_0$ of $\Sigma$-trees there clearly exists~$k$ such that the
cutting of all trees of~$M_0$ at level~$k$ yields trees pairwise
nonequivalent under~$\sim$---in other words, $j{\cdot}\partial_k{\cdot}m_0$
is a monomorphism. Choose $f\colon M_0\to M$ with $mf=j^*m_0$ to
conclude that, since
\[\overline{\partial}_k{\cdot}m{\cdot}f=j{\cdot}\partial_k{\cdot}m_0\]
is a monomorphism, also $\overline{\partial}_k{\cdot}m$~is a
monomorphism.

(2) The cone $(H\overline{\partial}_k)_{k\in\omega}$ is
collectively monomorphic. In fact, given distinct elements $x,y\in
HT_\Sigma$ there exists a finite nonempty set $m\colon M\to J^*$
with $x,y \in Hm[HM]$ , see 2.1~(ii). Choose~$k$ with
$\overline{\partial}_k{\cdot}m$ a monomorphism to conclude
$H\overline{\partial}_k(x)\neq H\overline{\partial}_k(y)$.

(3) The morphism $\varphi\colon HJ\to J$ of Definition~2.6 is a
monomorphism. In fact, this states that given two tuples
\[s=\sigma(s_1,\dots,s_n)\quad\text{and}\quad
r=\varrho(r_1,\dots,r_m)\]
in~$H_\Sigma I_{\Sigma^\perp}$ with $\sigma,\varrho\in\Sigma$, then
if $j\varphi_{\Sigma^\perp}u_{I_{\Sigma_\perp}}$~merges these tuples, then also
$\varepsilon_J{\cdot}H_\Sigma j$ merges them. In other words if $s$
and~$r$ are considered as trees in~$I_{\Sigma^\perp}$
(via~$\varphi_{\Sigma^\perp}$), then
\[s\sim r\quad\text{implies}\quad \varepsilon_J{\cdot}H_\Sigma
j\bigl(\sigma(s_i)\bigr) = \varepsilon_J{\cdot}H_\Sigma
j\bigl(\varrho(r_k)\bigr).\]
To prove this, we can restrict ourselves to the case that $s$~can be
obtained from~$r$ by a single  application of some
$\varepsilon$-equation. If the node of~$r$ at which this equation is
applied is not the root,
then that node lies in the subtree~$\varrho_{\llll}$
for some $\llll=0,\dots,m-\nolinebreak1$. This implies
\[s=\varrho(r_0,\dots,r_{\llll-1},s_{\llll},r_{\llll+1},\dots,r_{m-1})\]
and $j(s_{\llll})=j(r_{\llll})$, therefore
\[H_\Sigma j\bigl(\sigma(s_i)\bigr)=H_\Sigma
j\bigl(\varrho(r_k)\bigr).\]
If the above node is he root of~$r$, then the $\varepsilon$-equation
we apply has to have form
\[\sigma(x_0,\dots,x_{n-1})=\varrho(y_0,\dots,y_{m-1})\]
where $x_i$ and~$y_k$ are variables of a set~$X$, and for some
function $f\colon X\to I_{\Sigma^\perp}$ we have
\[s_i=f(x_i)\quad\text{and}\quad r_j=f(y_j).\]
Since the given equation is an $\varepsilon$-equation, i.e., the two
tuples are merged by~$\varepsilon_{X}$, the naturality
of~$\varepsilon$ yields
\begin{align*}
\varepsilon_J{\cdot}H_\Sigma j\bigl(\sigma(s_i)\bigr)
&=\varepsilon_J{\cdot}H_\Sigma j{\cdot}H_\Sigma f\bigl(\sigma(x_i)\bigr)\\
&=Hj{\cdot}Hf{\cdot}\varepsilon_X\bigl(\sigma(x_i)\bigr)\\
&=Hj{\cdot}Hf{\cdot}\varepsilon_X{\cdot}\bigl(\varrho(y_k)\bigr)\\
&=\varepsilon_J{\cdot}H_\Sigma j{\cdot}H\Sigma
f\bigl(\varrho(y_k)\bigr)\\
&=\varepsilon_J{\cdot}H_\Sigma j\bigl(\varrho(r_k)\bigr).
\end{align*}

(4) For each $k\in\omega$ the following squares
\vspace{1.5mm}
\[
\vcenter{
\xy
\xymatrixcolsep{3pc}
\xymatrix{
T_\Sigma
\ar[0,2]^-{\partial_{k+1}}
\ar[1,0]_{\psi_\Sigma}
&
&
I_{\Sigma^\perp}
\\
H_\Sigma T_\Sigma
\ar[0,1]_-{u_{T_\Sigma}}
&
H_{\Sigma^\perp}T_{\Sigma}
\ar[0,1]_-{H_{\Sigma^\perp}\partial_k}
&
H_{\Sigma^\perp}I_{\Sigma^\perp}
\ar[-1,0]_{\varphi_{\Sigma^\perp}}
}
\endxy
}
\qquad
\text{and}
\qquad
\vcenter{
\xy
\xymatrix{
J^*
\ar[0,1]^-{\overline{\partial}_{k+1}}
\ar[1,0]_{\varphi^*}
&
J
\\
HJ^*
\ar[0,1]_-{H\overline{\partial}_k}
&
HJ
\ar[-1,0]_{\varphi}
}
\endxy
}
\]\vspace{1mm}

\noindent
commute. In fact,the left-hand square commutes because for every
tree $s=\sigma(s_i)$ in~$T_k$ the cutting of~$s$ at
level~$k+\nolinebreak1$ can be performed by cutting each~$s_i$ at
level~$k$, i.e.,
\[\partial_{k+1}s=\sigma(\partial_ks_i).\]
To prove that the right-hand square commutes, use the fact that
$j^*\colon T_\Sigma\to J^*$ is an epimorphism with
\begin{align*}
\overline{\partial}_{k+1}j^*&=j{\cdot}\partial _{k+1}&&\text{by (1) above}\\
&=j{\cdot}\varphi_{\Sigma^\perp}{\cdot}H_{\Sigma^\perp}\partial_k{\cdot}
u_{T_\Sigma}{\cdot}\psi_\Sigma&&\text{see the left-hand square}\\
&=j{\cdot}\varphi_{\Sigma^\perp}{\cdot}u_{I_{\Sigma^\perp}}{\cdot}
H_\Sigma\partial_k{\cdot}\psi_{\Sigma}&&\text{by naturality of $u$}\\
&=\varphi{\cdot}\varepsilon_J{\cdot}H_\Sigma(j{\cdot}\partial_k){\cdot}\psi_\Sigma&&
\text{by 2.6}\\
&=\varphi{\cdot}H(j{\cdot}\partial_k){\cdot}\varepsilon_{T_\Sigma}{\cdot}
\psi_\Sigma&&\text{by naturality of $\varepsilon$}\\
&=\varphi{\cdot}H(\overline{\partial}_k{\cdot}j^*){\cdot} \varepsilon_{T_\Sigma}
{\cdot}\psi_\Sigma&&\text{by (1) above}\\
&=\varphi{\cdot}H\overline{\partial}_k{\cdot}\varphi^*{\cdot}j^*&&
\text{by Remark 2.11.}
\end{align*}

(5) The morphism $\psi^*\colon J^*\to HJ^*$ is well-defined in
the above Remark. In other words, the kernel equivalence of~$j^*$
contains the kernel equivalence of
\[\beta\overset{\text{def}}{=}Hj^*{\cdot}\varepsilon_{T_\Sigma}{\cdot}
\psi_\Sigma\colon T_\Sigma\to HJ^*.\]
Since the cone of all $\varphi{\cdot}H\overline{\partial}_k\colon
HJ^*\to J$ ($k\in\omega$) is monomorphic by (2) and~(3) above, it is
sufficient to show that the kernel equivalence of~$j^*$ contains
that of the cone $(\varphi{\cdot}H\overline{\partial}_k{\cdot}\beta)$. In
fact, we prove that
\[\varphi{\cdot}H\overline{\partial}_k{\cdot}\beta=j{\cdot}\partial_{k+1}\qquad
\text{for all $k\in\omega$,}\]
which proves the statement since (by definition of \sims\ and~$j^*$)
we have $j^*(s)=j^*(r)$ iff
$j{\cdot}\partial_{k+1}(s)=j{\cdot}\partial_{k+1}(r)$ for all $k\in\omega$.
The proof is as follows:
\begin{align*}
\varphi{\cdot}H\overline{\partial}_k{\cdot}\beta&=
\varphi{\cdot}(H\overline{\partial}_k{\cdot}\varepsilon_{J^*}){\cdot}
H_{\Sigma}j^*{\cdot}\psi_\Sigma&&
\text{by definition of $\beta$}\\
&=\varphi{\cdot}\varepsilon_J{\cdot}H_\Sigma(\overline{\partial}_kj^*)
{\cdot}\psi_\Sigma&&\text{by naturality of $\varepsilon$}\\
&=\varphi{\cdot}\varepsilon_J{\cdot}H_\Sigma(j{\cdot}\partial_k){\cdot}\psi_\Sigma&&
\text{by (1) above}\\
&=j{\cdot}\varphi_{\Sigma^\perp}{\cdot}u_{I_{\Sigma^\perp}}{\cdot}H_\Sigma
\partial_k{\cdot}\psi_\Sigma&&\text{by Definition 2.5}\\
&=j{\cdot}\varphi_{\Sigma^\perp}{\cdot}H_{\Sigma^\perp}\partial_k{\cdot}
u_{T_\Sigma}{\cdot}\psi_\Sigma&&\text{by naturality of $u$}\\
&=j{\cdot}\partial_{k+1}&&\text{by (4) above.}
\end{align*}

(6) The coalgebra~$J^*$ is terminal. In fact, given a coalgebra
$C\xrightarrow{\rule{1mm}{0mm}\gamma\rule{1mm}{0mm}}HC$, choose a
morphism
\[m\colon HC\to H_\Sigma C\qquad\text{with $\varepsilon_Cm=\id$.}\]
From the $H_\Sigma$-coalgebra
$C\xrightarrow{\rule{1mm}{0mm}\gamma\rule{1mm}{0mm}}
HC\xrightarrow{\rule{1mm}{0mm}m\rule{1mm}{0mm}}H_\Sigma C$ we have a
unique homomorphism $f\colon C\to T_\Sigma$ in~$\Coalg H_\Sigma$. We
obtain a homomorphism $j^*{\cdot}f\colon C\to J^*$ in~$\Coalg H$.
%\vspace{1.5mm}
%\[
%\psset{arrows=->,nodesep=3pt}
%\everypsbox{\scriptstyle}
%\begin{psmatrix}
%C&&&HC\\
%&HC&H_\Sigma C\\
%T_\Sigma&&H_\Sigma T_\Sigma\\
%&&&HT_\Sigma\\
%J^*&&&HJ^*
%\ncline{1,1}{1,4}^\gamma
%\ncline{1,1}{2,2}^\gamma
%\ncline{2,2}{2,3}^m
%\ncline{2,3}{1,4}^{\varepsilon_C}
%\ncline{1,1}{3,1}<f
%\ncline{2,3}{3,3}>{H_\Sigma f}
%\ncline{3,1}{3,3}^{\psi_\Sigma}
%\ncline{3,3}{4,4}^{\varepsilon_{T_\Sigma}}
%\ncline{3,1}{5,1}<{j^*}
%\ncline{4,4}{5,4}>{Hj^*}
%\ncline{5,1}{5,4}_{\varphi^*}
%\ncline{1,4}{4,4}>{Hf}
%\end{psmatrix}
%\]\vspace{1mm}

To prove the uniqueness of that homomorphism, let $h_1,h_2\colon
C\to J^*$ be homomorphisms of $H$-coalgebras. Since (1)~implies that
cone~$(\overline{\partial}_k)_{k\in\omega}$ is collectively
monomorphic, it is sufficient to prove
\[\overline{\partial}_kh_1=\overline{\partial}_kh_2\qquad
\text{for all $k\in\omega$.}\]
We proceed by induction in~$k$.

\textit{Case $k=0$}. The function~$\partial_0$ is constant (with
value~$\perp$), thus, it follows that $\overline{\partial}_0$~is
also constant.

\textit{Induction step}. Since $h_t$~is a homomorphism for $t=1,2$,
it follows from (4) above that
\[\overline{\partial}_{k+1}h_t=\varphi{\cdot}H \overline{\partial}_k
{\cdot}\varphi^*{\cdot}h_t=\varphi
{\cdot}H(\overline{\partial}_kh_t){\cdot}\gamma.\]
Thus, from $\overline{\partial}_kh_1=\overline{\partial}_kh_2$
we conclude
$\overline{\partial}_{k+1}h_1=\overline{\partial}_{k+1}h_2$.\qed
\end{pf}

\section{A Description of Completely Iterative Monads}

\begin{pv}{3.1}
We use the terminology of ~\cite{aamv}. Given a finitary endofunctor $H$,
presented as in 2.1 above, we denote by $\TT_\Sigma$ the free completely iterative monad on $H_\Sigma$ (of infinite $\Sigma$-trees) and by $\TT_H$ the free completely iterative monad on $H$. The signature obtained from $\Sigma$ by adding new nullary
operation symbols from a set $X$ is denoted by $\Sigma+X$.
\end{pv}

\begin{pv}{3.2 Remark}
For every finitary functor $H$, see 2.1, and every set $X$ we have a quotient
\[\varepsilon+\id_X\colon H_{\Sigma+X}=H_\Sigma({-})+X\to
H({-})+X.\]
Observe that $\varepsilon$-equations, as defined in
~2.1, are precisely the same as
$(\varepsilon+\id_X)$-equations. Denote by
\[\simx\]
the corresponding congruence on the
$(\Sigma+X)$-algebra~$I_{\Sigma^\perp+X}$, see Definition~2.6.
(For $X\neq\emptyset$ the choice of a new
symbol~$\perp$ is superfluous---instead, any element of~$X$ would
have done.)

Further, let
\[\simxs\]
denote the congruence on $T_\Sigma X=T_{\Sigma+X}$ of
Definition~2.9 given by applying the $\varepsilon$-equations
finitely or infinitely many times. That is, $\Sigma$-trees
$s$ and~$r$ over~$X$
are congruent iff $\partial_ks\simx\partial_kr$ holds
for all $k\in\omega$.
\end{pv}

\begin{pv}{3.3 Example}
For $H=\PP_{\!2}$ , see 2.2~(i), the congruence~\simx\ on
the algebra~$T_{\Sigma^\perp+X}$ (of all finite binary trees with
leaves labeled in $X+\{b,\perp\}$ and inner nodes labeled by~$\beta$)
is just the commutativity of the operation~$\beta$. And \simxs~is
the congruence on the algebra~$T_{\Sigma+X}$ (of all binary trees
with leaves labeled in $X+\{b\}$) which uses the commutativity
finitely or infinitely many times.
%
%\[
%\psset{levelsep=10mm,nodesep=3pt}
%\raisebox{19mm}{$
%\pstree{\TR{\beta}}{\TR{b}
%\pstree{\TR{\beta}}{\TR{x_1}
%\pstree{\TR{\beta}}{\TR{b}
%\pstree{\TR{\beta}}{\TR{x_2}
%\pstree{\TR{\beta}}{\TR{b}
%\Tp~[tnpos=b,tnsep=0pt]{\makebox[0pt]{%
%\kern5pt\raisebox{7pt}{.}\kern-.7pt\raisebox{3pt}{.}\kern-1.0pt%
%\raisebox{-.9pt}{.}}}}}}}}$}\quad\simxs\quad
%\raisebox{19mm}{$
%\pstree{\TR{\beta}}{%
%\pstree{\TR{\beta}}{%
%\pstree{\TR{\beta}}{%
%\pstree{\TR{\beta}}{%
%\pstree{\TR{\beta}}{%
%\Tp~[tnpos=b,tnsep=0pt]{\makebox[0pt]{%
%\kern-4.5pt\raisebox{-1.1pt}{.}\kern-.9pt\raisebox{3pt}{.}\kern-.9pt%
%\raisebox{7pt}{.}}}
%\TR{b}} \Tr{x_2}} \TR{b}} \TR{x_1}} \TR{b}}$}
%\]
\end{pv}

\begin{pv}{3.4 Theorem}
\itshape
For every finitary endofunctor~$H$ of~\Set\ a free completely
iterative monad~$\TT_H$ on~$H$ can be defined as the quotient of the
tree monad~$\TT_\Sigma$ modulo the monad congruence~\simxs\
\textnormal{(}$X$ a set\textnormal{)}.
\end{pv}

\begin{pv}{Remark}
More detailed: given a presentation as a quotient $\varepsilon\colon
H_\Sigma\to H$, then for the free completely iterative monad~$\TT_H$
on~$H$ we have a monad homomorphism $h\colon\TT_\Sigma\to\TT_H$ whose
components~$h_X$ are epimorphisms with the kernel equivalence~\simxs.
\end{pv}

\begin{pf}
The universal property of $\TT_\Sigma$ yields a monad homomorphism
\linebreak
$h\colon\TT_\Sigma\to\TT_H                        $ with surjective components. We know from
Theorem~2.11 that a terminal coalgebra, $TX$, of the finitary
endofunctor $H({-})+X$ can be obtained from the
tree coalgebra $T_\Sigma X=T_{\Sigma+X}$ modulo the
congruence~\simxs. Consequently, \simxs~is the kernel equivalence
of $h_X\colon T_\Sigma X\to TX$.\qed
\end{pf}

\section{Description of Iterative Monads}
\begin{pv}{4.1}
Recall from~\cite{be} or~\cite{e} that iterative monads are those ideal
monads $S$ which have a unique solution for every guarded equation
morphism $e\colon X\to S(Y+\nolinebreak Y)$ with $X$ and~$Y$ finite.
\end{pv}

\begin{pv}{Example}
(see ~\cite{ebt}). Recall from the Introduction that a tree in~$T_\Sigma X$ is called
\textit{rational} provided that it has only finitely subtrees.
The tree monad $\TT_\Sigma$ has a submonad~$\RR_\Sigma$
formed by all rational $\Sigma$-trees. More precisely, put
\[\RR_\Sigma=(R_\Sigma,\eta,\mu)\]
where $R_\Sigma$~assigns to every set~$X$ the subalgebra~$R_\Sigma
X$ of~$T_\Sigma X$ of all rational trees, $\eta_X$~sends~$x$ to the
singleton tree labeled by~$x$, and $\mu_X$~is given by tree tupling.
(It is easy to see that a tree-tupling of rational trees is always
rational.)
\end{pv}

\begin{pv}{4.2 Theorem}
(see ~\cite{ebt}).
\itshape
The rational-tree monad~$\RR_\Sigma$ is a free iterative monad on
the signature~$\Sigma$.
\end{pv}

\begin{pv}{4.3 Remark}
Given a finitary endofunctor~$H$ of\/~\Set, we define in~\cite{amv} a
subfunctor~$R$ of the above free iterative monad~$\TT_H$ on~$H$ by
\[RY=\bigcup e^\dag[Y]\]
where the union ranges over all flat equation morphisms $e\colon
X\to HX+\nolinebreak Y$ with $X$~finite and $e^dag$ denotes the unique
solution of  e ).
\end{pv}

\begin{pv}{4.4 Theorem}
(see AMV]).
\itshape
Every finitary endofunctor~$H$ of~\/~\Set\ generates a free
iterative monad, viz, the submonad~$\RR_H$ of~$\TT_H$ carried by the
above subfunctor~$R$.
\end{pv}

\begin{pv}{4.5 Notation}
Let $H$~be a finitary endofunctor of~\Set\ represented as a quotient
\[\varepsilon: H_\Sigma\to H.\]
For every set~$X$ we denote by~\vxs\ the congruence on the
rational-tree algebra~$R_\Sigma X$ which is the restriction of the
congruence~\simxs\ of 3.2. That is, two rational
$\Sigma$-trees $s$ and~$r$ over~$X$ are congruent iff $r$~can be
obtained from~$s$ by (potentially) infinite applications of
$\varepsilon$-equations---i.e., iff $\partial_ks\simx\partial_kr$
for all $k\in\nolinebreak\omega$.
\end{pv}

\begin{pv}{4.6 Theorem}
(Description of free iterative monads)
\itshape
For every finitary endofunctor~$H$ on\/~\Set\ a free iterative
monad~$\RR_H$ on~$H$ can be described as the quotient of the
rational-tree-monad~$\RR_\Sigma$ modulo the monad congruence~\vxs\
($X$ a set).
\end{pv}

\begin{pv}{Remark}
We thus exhibit, for every presentation of~$H$ as a quotient
$\varepsilon\colon H_\Sigma\to H$, a monad homomorphism
$h\colon\RR_\Sigma\to\RR$ whose components~$h_X$ are epimorphisms
with the kernel equivalence~\vxs.
\end{pv}

\begin{pf}
Denote by
\[k\colon\TT_\Sigma\to\TT_H\]
the quotient of Theorem~3.4. Observe that for every flat equation
morphism $e\colon X\to H_\Sigma X+\nolinebreak Y$ w.r.t~$\Sigma$ we
obtain a flat equation morphism
\[\overline{e}\equiv X\xrightarrow{\rule{1mm}{0mm}e\rule{1mm}{0mm}}
H_\Sigma X+Y
\xrightarrow{\rule{1mm}{0mm}\varepsilon_X+\id_Y\rule{1mm}{0mm}}
HX+Y\]
w.r.t~$H$ such that the following triangle
\begin{equation}
\vcenter{
\xy
\xymatrix{
&
X
\ar[1,-1]_{e^\dag}
\ar[1,1]^{\overline{e}^\dag}
&
\\
T_\Sigma Y
\ar[0,2]_-{k_Y}
&
&
TY
}
\endxy
}
\end{equation}
commutes. In fact, since $e^\dag$~is a coalgebra homomorphism w.r.t
$H_\Sigma({-})+Y$ (see ~\cite{aamv}, Solution Lemma 3.4 ) so
is~$k_Ye^\dag$, and this clearly implies that $k_Ye^\dag$~is a
coalgebra homomorphism w.r.t
$H_\Sigma({-})+\nolinebreak Y$. Since $TY$~is
terminal, we conclude $\overline{e}^\dag=k_Ye^\dag$.

This implies that $k_Y$~has a domain-codomain restriction
$h_Y\colon R_\Sigma Y\to RY$. In fact, every element~$r$ of~$R_\Sigma
Y$ is a solution of some flat system $e\colon X\to H_\Sigma
X+\nolinebreak Y$ with $X$~finite, more precisely, $r=e^\dag(x)$ for
some $x\in X$. Then
\[k_Y(r)=k_Ye^\dag(x)=\overline{e}(x)\in RY.\]
Since $k\colon T_\Sigma\to T$ is a natural transformation, it
follows that the maps~$h_Y$ form a natural transformation $h\colon
R_\Sigma\to R$. Moreover, $h$~is a monad morphism,
$h\colon\RR_\Sigma\to\RR$, because it is a restriction of the monad
morphism $k\colon\TT_\Sigma\to\TT$ to submonads.

It remains to prove that $h_Y$~is surjective. To this end, for every
flat equation $\overline{e}\colon X\to HX+\nolinebreak Y$ choose a
splitting of~$\varepsilon_X$:
\[u\colon HX\to H_\Sigma X,\qquad \varepsilon_Xu=\id\]
and consider the flat equation $e=(u+\id_Y)\overline{e}\colon X\to
H_\Sigma X+\nolinebreak Y$ w.r.t.~$H_\Sigma$. Then the above
triangle~(3) commutes. To verify this, we only need to prove
that $e^\dag$~is a coalgebra homomorphism w.r.t.
$H({-})+Y$ from~$\overline{e}$ to the coalgebra
$T_\Sigma Y\cong H_\Sigma T_\Sigma
Y+Y\xrightarrow{\rule{1mm}{0mm}\varepsilon_{T_\Sigma
Y}+\id\rule{1mm}{0mm}}H(T_\Sigma Y)+\nolinebreak Y$. In fact, from
the Solution Lemma of ~\cite{aamv} we know that $e^\dag$~is a coalgebra homomorphism
from~$e$ to~$T_\Sigma Y$. That is, in the following diagram
\vspace{1mm}

\[
\xy
\xymatrix{
X
\ar[0,1]^-{\overline{e}}
\ar[3,0]_{e^\dag}
&
HX+Y
\ar[0,2]^-{u+\id_Y}
\ar @{=} [1,1]
&
&
H_\Sigma X+Y
\ar[3,0]^{H_\Sigma e^\dag+\id_Y}
\ar[1,-1]_{\varepsilon_X+\id_Y\phantom{MM}}
\\
&
&
HX+Y
\ar[1,0]^{He^\dag+\id_Y}
&
\\
&
&
H(T_\Sigma Y)+Y
&
\\
T_\Sigma Y
\ar[0,3]_-{\cong}
&
&
&
H_\Sigma(T_\Sigma Y)+Y
\ar[-1,-1]^{\varepsilon_{T_\Sigma Y}+\id_Y\phantom{MM}}
}
\endxy
\]
\vspace{1mm}

\noindent
the outward square commutes. Since the right-hand square commutes by
naturality of~$\varepsilon$, it follows that $e^\dag$~is a
homomorphism from~$\overline{e}$ w.r.t $H({-})+Y$, as
requested. This shows that $h_Y$~is surjective: every element
$\overline{r}\in RY$ has the form
$\overline{r}=\overline{e}^\dag(x)$ for some flat equation
$\overline{e}\colon X\to HX+\nolinebreak Y$ with $X$~finite and some
$x\in X$, and then we have
\[\overline{r}=h_Y\bigl(e^\dag(x)\bigr)\qquad\text{with}\qquad
e^\dag(x)\in R_\Sigma Y.\]
\mbox{}\vspace{-34pt}\par\mbox{}\qed
\end{pf}



\begin{thebibliography}{\text{MMMM}}
\bibitem[\text{AAMV}]{aamv}
Aczel,~P., Ad\'amek,~J., Milius,~S. and Ve\-le\-bil,~J.: ``Infinite Trees and
Completely Iterative Theories''; to appear in Theor. Comps. Science.

%\bibitem[\text{A}]{a}
%Ad\'amek,~J.: ``Free Algebras and Automata Realizations in the
%Language of Categoies''; Comment. Math. Univ. Carolin\ae, 15(1974),
%\mbox{589--602}.

\bibitem[\text{AMV}]{amv}
Ad\'amek,~J., Milius,~S. and Velebil,~J.: ``On Rational Monads and Free Iterative Theories''; Electr. Notes Theor. Comp. Science 69 (2002)

\bibitem[\text{AT}]{at}
Ad\'amek,~J. and Trnkov\'a,~V.: ``Automata and Algebras in
Categories''; Academic Publishers, Dordrecht (1990)

\bibitem[\text{B}]{b}
Barr,~M.: ``Terminal Coalgebras in Well-founded Set Theory''; Theor.
Comp. Science 124(1994), \mbox{182--192}.

\bibitem[\text{BE}]{be}
Bloom, S.~L. and \'Esik,~Z.: ``Iteration Theories: The Equational
Logic of Iterative Processes''; EATCS Monograph Series on
Theoretical Computer Science, Springer-Verlag (1993)

\bibitem[\text{E}]{e}
Elgot, C.~C.: ``Monadic Computation and Iterative Algebraic
Theories''; in: Logic Colloquium~'73 (eds: H.~E. Rose and J.~C.
Shepherdson), North-Holland Publishers, Amsterodam (1975),
\mbox{175--230}.

\bibitem[\text{EBT}]{ebt}
Elgot, C.~C., Bloom, S.~L. and Tindell,~R.: ``On the Algebraic
Structure of Rooted Trees''; J. Comp. Syst. Sciences, 16(1978),
\mbox{361--399}.

\bibitem[\text{L}]{l}
Lambek,~J.: ``A Fixpoint Theorem for Complete Categories''; Math.
Z., 103(1968), \mbox{151--161}.

\bibitem[\text{M}]{m}
Moss,~L.: ``Parametric Corecursion''; Theoretical Computer Science,
260(2001), \mbox{139--163}.

\end{thebibliography}


\end{document}

}
